\documentclass[10pt]{article}
\usepackage[letterpaper,margin=0.68in]{geometry}
\usepackage{amsmath,amssymb}
\usepackage{enumitem}
\usepackage{xcolor}
\usepackage{fancyhdr}
\usepackage{microtype}
\usepackage[most]{tcolorbox}
\definecolor{olive}{HTML}{4D5430}
\definecolor{gold}{HTML}{9A7B22}
\definecolor{pale}{HTML}{F5F1DC}
\definecolor{softgreen}{HTML}{EDF0E2}
\pagestyle{fancy}\fancyhf{}\lhead{\textcolor{olive}{MATH\& 107 Answer Key}}\rhead{\textcolor{gold}{Fall 2026}}\cfoot{\thepage}
\setlength{\headheight}{14pt}
\setlength{\parindent}{0pt}\setlength{\parskip}{5pt}
\setlist[enumerate]{leftmargin=*,itemsep=5pt,topsep=4pt}
\newtcolorbox{solbox}[1]{colback=pale,colframe=olive,boxrule=.7pt,arc=2mm,title=\textbf{#1},fonttitle=\color{olive},left=2mm,right=2mm,top=1.5mm,bottom=1.5mm}
\newcommand{\modulekey}[2]{\clearpage{\Large\bfseries\textcolor{olive}{Module #1: #2 — Answer Key}}\par\textcolor{gold}{\rule{\linewidth}{1.2pt}}\par}
\begin{document}
\modulekey{7}{Cost, Revenue, and Profit}
\textbf{Prequiz}
\begin{enumerate}
\item \(F=180\), \(v=6\), \(p=18\). Thus
\[
C(q)=\boxed{180+6q},\qquad R(q)=\boxed{18q},\qquad \pi(q)=\boxed{12q-180}.
\]
\item
\[
C(25)=180+6(25)=\boxed{330},
\]
\[
R(25)=18(25)=\boxed{450},
\]
\[
\pi(25)=450-330=\boxed{120}.
\]
The business has a profit.
\item
\[
q^*=\frac{180}{18-6}=\boxed{15}.
\]
At \(q=14\), profit is \(-12\); at \(q=16\), profit is \(+12\).
\item Marginal cost \(=\boxed4\), marginal revenue \(=\boxed9\), marginal profit \(=\boxed5\). Break-even:
\[
q^*=\frac{250}{9-4}=\boxed{50}.
\]
Ten additional meals after break-even add \(10(5)=\boxed{\$50}\) profit.
\item To the right of break-even, \textbf{revenue is larger than cost}, so profit is positive.
\end{enumerate}

\textbf{Matching quiz}
\begin{enumerate}
\item
\[
C(q)=\boxed{240+3q},\qquad R(q)=\boxed{9q},\qquad \pi(q)=\boxed{6q-240}.
\]
\item At \(q=50\):
\[
C=\boxed{390},\qquad R=\boxed{450},\qquad \pi=\boxed{60}.
\]
\item
\[
q^*=\frac{240}{9-3}=\boxed{40}.
\]
The minimum whole number for nonnegative profit is \(40\).
\item Marginal profit \(=15-7=\boxed8\). Twelve additional units add \(12(8)=\boxed{\$96}\) profit.
\item Revenue is higher than cost, so profit is positive.
\end{enumerate}
\end{document}
